\documentclass{ctexart}
\usepackage{amsmath, amssymb}
\usepackage{geometry}
\geometry{a4paper, margin=2cm}

\title{\textbf{第5章 质点的动量矩与动量矩守恒} --- 例题}
\author{}
\date{}

\begin{document}
\maketitle

\begin{enumerate}

\item（例 6.10）质量为 $m$ 的小球，悬挂在一根轻杆下端，杆的上端固定在一个可绕竖直轴旋转的铰链上。设小球在水平面内作半径为 $r$ 的匀速率圆周运动，角速度为 $\omega$，轻杆与竖直方向之间的夹角为 $\alpha$。求小球对 $A$ 点（铰链）和 $B$ 点（杆与小球连接点）的动量矩。

\item 一质点 $m$，速度为 $v$，$A$、$B$、$C$ 分别为三个参考点，此时 $m$ 相对三个点的距离分别为 $d_1$、$d_2$、$d_3$。求此时刻质点对三个参考点的动量矩。

\item（例题 1）系在绳子的一端，绳的另一端穿过水平光滑桌面中央的小洞 $O$，起初下面用手拉着不动，质点在桌面上绕 $O$ 作匀速圆周运动。然后慢慢下拉绳子，使它在桌面上那一段缩短。质点绕 $O$ 的角速度 $\omega$ 如何随半径 $r$ 变化？

\item（例题 2）两个同样重的小孩，各抓着跨过滑轮绳子的两端。一个孩子用力向上爬，另一个则抓住绳子不动。若滑轮的质量和轴上的摩擦都可忽略，哪一个小孩先到达滑轮？又：两个小孩重量不等时情况如何？

\item（书中例题 6.11，p.214）人造卫星在椭圆轨道上运行，地球中心可看作固定点，近地点离地面的距离为 $439\,\text{km}$，远地点离地面的距离为 $2384\,\text{km}$，近地点速度为 $8.12\,\text{km/s}$，地球半径为 $6370\,\text{km}$。求卫星在远地点的速度。

\item（书中例题 6.12，p.215，重点）质量为 $m$ 的小球系在绳子的一端，绳穿过一铅直套管，使小球限制在一光滑水平面上运动。先使小球以速度 $v_0$ 绕管心作半径为 $r_0$ 的圆周运动，然后向下拉绳，使小球轨迹最后成为半径为 $r$ 的圆。试求：
\begin{enumerate}
  \item 小球距管心 $r$ 时速度 $v$ 的大小；
  \item 绳从 $r_0$ 缩短到 $r$ 过程中，力 $F$ 所作的功。
\end{enumerate}

\item 发射一宇宙飞船去考察一质量为 $M$、半径为 $R$ 的行星，当飞船静止于空间距行星中心 $4R$ 时，以速度 $v_0$ 发射一质量为 $m$ 的仪器。要使该仪器恰好掠过行星表面，求 $\theta$ 角及着陆滑行的初速度。

\end{enumerate}

\end{document}
